\documentclass{book}
\usepackage{amsmath}
\usepackage{amsthm}

\theoremstyle{definition}
\newtheorem{definition}{Definition}[section]
\newtheorem{proposition}[definition]{Proposition}

\setcounter{chapter}{1}

\begin{document}

\section{Definition}

\begin{definition}
Let $F$ be a field. A \textbf{vector space} over $F$ is a set $V$ equipped with two operations
\[
\begin{aligned}
+       &\colon V \times V \to V, \quad (u, v) \mapsto u + v, \\
\cdot   &\colon F \times V \to V, \quad (\lambda, v) \mapsto \lambda v,
\end{aligned}
\]
satisfying the following axioms:

\begin{enumerate}
\item
    $(V, +)$ is an abelian group. That is, for all $u, v, w \in V$,
    \begin{itemize}
        \item $(u + v) + w = u + (v + w)$,
        \item there exists $0 \in V$ such that $v + 0 = v$,
        \item for each $v \in V$ there exists $-v \in V$ such that $v + (-v) = 0$,
        \item $u + v = v + u$.
    \end{itemize}
\item
    Scalar multiplication is compatible with field operations: for all $\lambda, \mu \in F$ and $u, v \in V$,
    \begin{itemize}
        \item $\lambda(u + v) = \lambda u + \lambda v$,
        \item $(\lambda + \mu)v = \lambda v + \mu v$,
        \item $(\lambda\mu)v = \lambda(\mu v)$,
        \item $1v = v$, where $1$ is the multiplicative identity of $F$.
    \end{itemize}
\end{enumerate}

Elements of $V$ are called \textbf{vectors}; elements of $F$ are called \textbf{scalars}.
\end{definition}

\begin{definition}
Let $S$ be any set and $F$ a field. The set
\[
F^{S} = \{ f \mid f \colon S \to F \}
\]
of all functions from $S$ to $F$ is a vector space over $F$ with operations defined pointwise:
\[
\begin{aligned}
(f + g)(x) &= f(x) + g(x), \\
(\lambda f)(x) &= \lambda \cdot f(x),
\end{aligned}
\]
for all $f, g \in F^{S}$, $\lambda \in F$, and $x \in S$.
\end{definition}

\begin{proposition}
$F^{S}$ is a vector space over $F$.
\end{proposition}

\begin{proof}
Left as an exercise.
\end{proof}

\end{document}
